\documentclass{article}
\title{\textbf{\huge \underline{Experts}}}
\date{}
\usepackage{amsmath,amssymb}
\usepackage{amsthm}
\usepackage{xcolor}
\newtheorem{theorem}{Theorem}[section]
\newtheorem{defination}{Defination}[section]
\begin{document}
	\maketitle
	\section{Analytic Equations}
	The famous formula of Leibniz involving $\pi$ is
	\begin{equation}
		\label{eq:1}
		\frac{\pi}{4} = 1-\frac{1}{3}+\frac{1}{5}-\frac{1}{7}+\frac{1}{9}-\cdots=\sum_{k\ge 0}^{} \frac{(-1)^k}{2k+1}.
	\end{equation}
	However, \eqref{eq:1} converges far too slowly to be a practical method of approximating $\pi$.
	
	Euler’s integral formula for $n!$ is
	\begin{equation*}
		n!=\int_{0}^{\infty} x^n e^{-x} dx.
	\end{equation*}
	The right side makes sense at $n = 1/2$ and can be evaluated explicitly:
	\begin{equation*}
		\int_{0}^{\infty} \sqrt{x} e^{-x} dx=\frac{\sqrt{\pi}}{2}.
	\end{equation*}
	This suggests defining $(1/2)! = \frac{\sqrt{\pi}}{2}$.
	
	For $a \in R$ and a function $f : \mathbb{R} \rightarrow \mathbb{R}$, the derivative of $f$ at $a$ is
	\begin{align*}
		f'(a) &= \lim_{h\rightarrow0} \frac{f(a+h)-f(a)}{h} \\
		&= \lim_{x\rightarrow a} \frac{f(x)-f(a)}{x-a}
	\end{align*}
	More formulas like this can be found in the book \cite[\textcolor{red}{Rud76}]{Ru76}.
	\section{Areas}
	If we are dealing with lattice polygons that are not simple, then we can use the following
	extension of Pick's Theorem.
	\begin{theorem}
		\label{th:1}
			For a lattice polygon $P$ containing $h$ holes with a boundary of $h + 1$ simple closed curves, its area is given by
			\[
				A(P)=I(P)+\frac{B(P)}{2}+h-1,
			\]
			where $I(P)$ is the number of interior lattice points of $P$ and $B(P)$ is the number of boundary
			lattice points of $P$.
	\end{theorem}

	We now review the definition of Euler characteristic for a convex polyhedron.
	\begin{defination}
		The Euler characteristic of a convex polyhedron $P$ with $V$ vertices, $E$
		edges, and $F$ faces is $\chi(P) = V - E + F.$
	\end{defination}


	Next, we recall Euler’s formula.
	
	\begin{theorem}
		(Euler`s Formula). For any convex polygon $P$ with Euler characteristic $\chi(P)$,
		we have that
		\begin{equation}
			\label{eq:2}
			\chi(P)=2.
		\end{equation}
	\end{theorem}


We will now generalize Theorem \ref{th:1} by incorporating the Euler characteristic, which will
allow us to find the area $A(P)$ of any lattice polygon $P$ that is the union of a finite number
of simple lattice polygons. This more general area formula is given by
\[
\begin{pmatrix}
	a & b \\
	c & d
\end{pmatrix}
\]


\end{document}